\documentclass[../../../include/open-logic-section]{subfiles}

\begin{document}

\olfileid{sth}{spine}{valpha}

\olsection{Netepake Tataran Minangka $V_\alpha$}

Ing \olref[sth][ordinals][]{chap}, kita wis netepake
urutan becik lan ordinal (miturut von Neumann).
Ing bab iki, kita bakal nggunakake loro gagasan iku
kanggo nggambarake hierarki himpunan \emph{dhewe}.
Kanggo iku, elinga yen ing
\olref[sth][ordinals][opps]{sec} kita wis netepake
ordinal penerus lan ordinal limit. Gagasan-gagasan
iki digunakake ing definisi sabanjure:

\begin{defn}\ollabel{defValphas}
	\begin{align*}
	V_\emptyset &\defis \emptyset\\
	V_{\ordsucc{\alpha}} &\defis \Pow{V_\alpha} & & 
	\text{kanggo sembarang ordinal }\alpha\\
	V_{\alpha} &\defis \bigcup_{\gamma < \alpha} V_\gamma & & 
	\text{yen }\alpha\text{ ordinal limit}
\end{align*}
\end{defn}
\noindent
Iki bakal dadi definisi nganggo \emph{rekursi
transfinit} ing ordinal. Ing bab iki, kita bisa
mbandhingake karo definisi rekursif fungsi ing
wilangan asli.\footnote{Bandhingna\ definisi
panjumlahan, pingan, lan pemangkatan ing
\olref[sfr][infinite][dedekind]{sec}.} Kaya ing
wilangan asli, kita netepake kasus dhasar lan kasus
penerus; nanging kanggo ordinal, kita uga kudu
nerangake kelakuan kasus \emph{limit}.

Definisi himpunan-himpunan $V_\alpha$ iki bakal dadi
tonggak kang wigati. Sadurunge, kita nerangake
hierarki himpunan kanthi cara informal: himpunan
kawangun miturut \emph{tataran}. Himpunan-himpunan
$V_\alpha$ iku sejatine tataran-tataran mau.
Nanging iki minangka cara nggambarake tataran kanthi
\emph{internal}. Kita ora mung ngusulake sawijining
\emph{model} kanggo teori iki; kita
netepake tataran \emph{ing sajroning} teori himpunan
kita dhewe.

\end{document}
